\rtmtight
\begin{tblr}{width=\textwidth,colspec={|Q[c,m,2.1cm]|X[l,m]|},hlines,vlines,hspan=minimal,rowsep=1.5pt,colsep=3pt,row{1}={bg=headgray,font=\bfseries}}
\SetCell{c,m}\hfil\logo\hfil & \textbf{POLITEKNIK NEGERI MALANG}\par\textbf{JURUSAN TEKNOLOGI INFORMASI}\par\textbf{PROGRAM STUDI: D4 TEKNIK INFORMATIKA} \\
\SetCell[c=2]{l} \textbf{RENCANA TUGAS MAHASISWA (RTM) DAN RUBRIK PENILAIAN} & \\
Nama Mata Kuliah & Magang \\
Kode Mata Kuliah & RTI257001 \quad \textbf{sks / Jam perminggu:} 20 SKS / 40 jam/minggu \quad \textbf{Semester:} 7 \\
Dosen Pengampu & Tim Pengajar Program Studi D4 TI \\
\end{tblr}
\assessmentblock{Logbook Kegiatan Magang}{Portofolio dan Case Method}{SCPMK0301-05701, SCPMK0804-05703}{Mahasiswa mengisi logbook kegiatan harian selama magang yang mencatat aktivitas kerja, pencapaian, dan refleksi profesional secara berkala.}{Mengisi logbook harian, mendokumentasikan aktivitas dan hasil kerja, melakukan refleksi mingguan, dan mendapatkan validasi pembimbing lapangan.}{Logbook digital/fisik yang terisi lengkap, divalidasi pembimbing lapangan, dan laporan ringkas kemajuan berkala.}{Kelengkapan dan konsistensi pengisian logbook & 8 & Logbook terisi setiap hari kerja, mencakup aktivitas, waktu, hasil, dan refleksi; divalidasi pembimbing lapangan secara berkala. & Logbook terisi sebagian besar hari kerja dengan informasi yang cukup, tetapi beberapa entri tidak lengkap atau tidak tervalidasi. & Logbook banyak hari kosong, isian tidak informatif, atau tidak divalidasi pembimbing lapangan. \\ \\
Kualitas refleksi dan dokumentasi aktivitas kerja & 7 & Refleksi menunjukkan pemahaman mendalam terhadap tugas industri, menghubungkan pengalaman dengan kompetensi akademis, dan mengidentifikasi pembelajaran. & Refleksi cukup informatif tetapi hubungan dengan kompetensi akademis atau identifikasi pembelajaran masih terbatas. & Refleksi bersifat deskriptif semata tanpa analisis atau koneksi dengan pembelajaran akademis. \\ \\
Ketepatan waktu dan kemandirian selama magang & 5 & Semua target logbook dipenuhi tepat waktu, mahasiswa menunjukkan inisiatif tinggi dan kemandirian dalam bekerja. & Sebagian besar target dipenuhi tepat waktu dengan sedikit keterlambatan atau kebutuhan bimbingan tambahan. & Banyak target terlambat atau mahasiswa terlalu bergantung pada pembimbing untuk mengisi logbook. \\}

\assessmentblock{Penilaian Pembimbing Lapangan}{Penilaian Industri}{SCPMK0804-05703, SCPMK0607-05704}{Penilaian dari pembimbing industri terhadap kinerja mahasiswa selama magang mencakup profesionalisme, penyelesaian tugas, dan penerapan kompetensi rekayasa perangkat lunak.}{Pembimbing lapangan mengobservasi kinerja harian, mengevaluasi hasil kerja, memberikan penilaian mid-term dan akhir menggunakan formulir penilaian standar.}{Formulir penilaian pembimbing lapangan mid-term dan akhir yang telah diisi dan ditandatangani.}{Profesionalisme, disiplin, dan etika kerja & 8 & Mahasiswa konsisten hadir tepat waktu, berpakaian profesional, berkomunikasi dengan santun, dan mematuhi aturan perusahaan tanpa pengingat. & Profesionalisme umumnya baik tetapi kadang ada keterlambatan, pelanggaran kecil aturan, atau komunikasi yang kurang tepat. & Sering terlambat, tidak mematuhi aturan perusahaan, atau menunjukkan perilaku tidak profesional berulang kali. \\ \\
Kualitas dan penyelesaian tugas industri & 12 & Seluruh tugas yang diberikan diselesaikan tepat waktu dengan kualitas memuaskan atau melebihi ekspektasi pembimbing lapangan. & Sebagian besar tugas diselesaikan tepat waktu dengan kualitas yang memadai; beberapa memerlukan revisi. & Banyak tugas tidak selesai tepat waktu, kualitas di bawah standar, atau memerlukan koreksi besar. \\ \\
Penerapan kompetensi rekayasa perangkat lunak & 10 & Mahasiswa secara nyata menerapkan kompetensi RPL (desain, implementasi, pengujian, atau dokumentasi) dalam tugas industri dan memberikan kontribusi teknis yang diapresiasi. & Kompetensi RPL diterapkan sebagian dalam tugas; kontribusi teknis ada tetapi masih terbatas pada tugas sederhana. & Kompetensi RPL tidak terlihat dalam pekerjaan, atau mahasiswa tidak berkontribusi secara teknis. \\}

\assessmentblock{Laporan Magang}{Case Method dan penulisan ilmiah}{SCPMK0301-05701, SCPMK0303-05702}{Mahasiswa menyusun laporan magang yang komprehensif mencakup latar belakang perusahaan, deskripsi pekerjaan, analisis pengalaman, dan refleksi kompetensi yang diaplikasikan.}{Mengumpulkan data dan dokumentasi selama magang, menyusun laporan sesuai sistematika Polinema, bimbingan dengan pembimbing akademik, dan revisi berdasarkan masukan.}{Naskah laporan magang lengkap sesuai pedoman Polinema, dilengkapi lampiran pendukung (surat keterangan magang, logbook, dokumentasi).}{Kelengkapan struktur dan konten laporan & 7 & Laporan memuat semua bab yang dipersyaratkan (pendahuluan, profil perusahaan, pelaksanaan magang, penutup) dengan konten yang lengkap, akurat, dan relevan. & Struktur laporan lengkap tetapi beberapa bab memiliki konten yang kurang mendalam atau kurang akurat. & Struktur laporan tidak lengkap, bab penting hilang, atau konten tidak mencerminkan pengalaman magang yang sesungguhnya. \\ \\
Kualitas analisis dan pembahasan pengalaman kerja & 8 & Analisis menghubungkan pengalaman industri dengan teori akademis secara kritis, mengidentifikasi kompetensi yang berkembang, dan memberikan insight berharga. & Analisis cukup baik tetapi koneksi teori-praktik atau kedalaman refleksi masih terbatas. & Laporan bersifat deskriptif tanpa analisis, atau tidak ada keterkaitan yang jelas antara pengalaman dan kompetensi akademis. \\ \\
Tata tulis, referensi, dan penyajian & 5 & Bahasa baku, tata tulis sesuai pedoman, referensi terdokumentasi dengan benar, dan laporan tersaji rapi serta profesional. & Tata tulis umumnya baik tetapi ada kesalahan penulisan, format referensi yang tidak konsisten, atau penyajian yang kurang rapi. & Banyak kesalahan bahasa, format tidak sesuai pedoman, atau referensi tidak terdokumentasi. \\}

\assessmentblock{Sidang Magang}{Presentasi dan defense}{SCPMK0303-05702, SCPMK0607-05704}{Mahasiswa mempresentasikan dan mempertahankan laporan magang di hadapan tim penguji (pembimbing akademik dan penguji internal), menunjukkan penguasaan atas pengalaman dan kompetensi industri yang diperoleh.}{Menyiapkan bahan presentasi, memaparkan ringkasan magang dan kompetensi yang dicapai, menjawab pertanyaan penguji secara analitis.}{Presentasi sidang (slide), naskah laporan akhir yang telah direvisi, dan berita acara sidang.}{Kualitas presentasi dan komunikasi lisan & 10 & Presentasi terstruktur, disampaikan dengan percaya diri dalam bahasa yang lugas, durasi terkontrol, dan materi mudah dipahami penguji. & Presentasi cukup terstruktur dan dapat dipahami, tetapi penyampaian kurang percaya diri atau durasi tidak terkontrol. & Presentasi tidak terstruktur, bahasa tidak profesional, atau materi sulit dipahami penguji. \\ \\
Kemampuan menjawab pertanyaan penguji & 10 & Semua pertanyaan dijawab dengan tepat, didukung argumen logis dan bukti dari pengalaman lapangan, menunjukkan pemahaman mendalam. & Sebagian besar pertanyaan dijawab dengan cukup baik, tetapi beberapa jawaban kurang didukung bukti atau argumen. & Banyak pertanyaan tidak dapat dijawab, jawaban tidak relevan, atau tidak dapat mempertahankan konten laporan. \\ \\
Keterkaitan kompetensi industri dengan akademis & 10 & Mahasiswa dapat menjelaskan secara eksplisit kompetensi RPL yang diterapkan di industri, menghubungkan dengan teori akademis, dan merumuskan nilai tambah pengalaman magang. & Keterkaitan kompetensi industri dan akademis dijelaskan tetapi masih umum atau kurang spesifik. & Tidak dapat menjelaskan keterkaitan antara pengalaman industri dan kompetensi akademis, atau pembahasan sangat dangkal. \\}

\begin{tblr}{width=\textwidth,colspec={|Q[l,m,3.2cm]|X[l,m]|},hlines,vlines,hspan=minimal,row{1-3}={bg=headgray},rowsep=1.5pt,colsep=3pt}
Jadwal Pelaksanaan: & Minggu 1--17 sesuai RPS. Setiap evaluasi memiliki RTM dan rubrik multi-kriteria. \\
Lain-Lain Yang Diperlukan: & LMS, logbook digital, perangkat kerja industri, dan perangkat presentasi. \\
Pustaka: & Mengacu pustaka utama dan pendukung pada bagian RPS. \\
\end{tblr}
